% =============================================================================
% preambulo/comandos.tex  --  Comandos auxiliares e macros do autor
% =============================================================================

% ----- TODO destacado no PDF (sem dependência do pacote todonotes) -----
% Uso: \todoex{texto da pendência}
\newcommand{\todoex}[1]{%
  \textcolor{red!70!black}{\textbf{[\,TODO:\,#1\,]}}%
}

% ----- Marcação inline para revisão -----
\newcommand{\rev}[1]{\textcolor{blue!60!black}{\textit{#1}}}

% ----- OBS para elementos fictícios (resultados preliminares a substituir) -----
% Uso: \obsfic{breve nota}  --> aparece em vermelho, em itálico, com prefixo OBS.
\newcommand{\obsfic}[1]{\textcolor{red!75!black}{\textit{\small OBS: #1}}}

% ----- Termos técnicos em inglês (itálico) -----
\newcommand{\estr}[1]{\textit{#1}}

% ----- Atalhos para nomes de modelos/datasets (manter consistência) -----
\newcommand{\modeloMD}{MegaDetector}
\newcommand{\modeloSN}{SpeciesNet}
\newcommand{\modeloCatFLW}{CatFLW}
\newcommand{\modeloPetFace}{PetFace}
\newcommand{\modeloWildlife}{WildlifeReID-10k}

% ----- Métricas (para garantir grafia consistente) -----
\newcommand{\mAP}{\mbox{mAP}}
\newcommand{\Fscore}{F\textsubscript{1}}
\newcommand{\RankK}[1]{Rank-#1}

% ----- Compatibilidade com .bbl gerado pelo abntex2-alf-USPSC.bst -----
% O estilo abntex2-alf emite comandos legados de LaTeX 2.09 (\tt, \it, \bf,
% \sf, \sc, \rm, \sl) e do pacote html (\htmladdnormallink). A classe memoir
% bloqueia esses comandos. Redefinimos para equivalentes modernos.
\let\tt\ttfamily
\let\it\itshape
\let\bf\bfseries
\let\sf\sffamily
\let\sc\scshape
\let\rm\rmfamily
\let\sl\slshape
% Stub para \htmladdnormallink{texto}{url}; em PDF imprime só o texto.
\providecommand{\htmladdnormallink}[2]{#1}
\providecommand{\htmladdnormallinkfoot}[2]{#1}

\raggedbottom